\documentclass[11pt]{article}
\usepackage[margin=1.15in]{geometry}
\usepackage{amsmath,amssymb,amsthm}
\usepackage{hyperref}

\theoremstyle{plain}
\newtheorem{theorem}{Theorem}
\newtheorem{lemma}[theorem]{Lemma}
\newtheorem{proposition}[theorem]{Proposition}
\newtheorem{corollary}[theorem]{Corollary}
\newtheorem{conjecture}[theorem]{Conjecture}
\theoremstyle{definition}
\newtheorem{remark}[theorem]{Remark}
\newtheorem{question}[theorem]{Question}

\newcommand{\Cfour}{C_4}
\newcommand{\Ceight}{C_8}
\newcommand{\Csixteen}{C_{16}}

\title{Powers of two as cycle lengths in graphs embedded in surfaces}
\author{}
\date{}

\begin{document}
\maketitle

\begin{abstract}
The Erd\H{o}s--Gy\'arf\'as conjecture asserts that every finite simple graph of
minimum degree at least $3$ contains a cycle whose length is a power of two. We
prove two new cases. First, every planar graph of minimum degree at least $4$
contains a cycle of length $4$ or $8$; the same holds for graphs embeddable in
the projective plane, and for $2$-connected graphs embeddable in the torus or
the Klein bottle. Second, every bipartite graph of minimum degree at least $3$
embeddable in the torus contains a cycle of length $4$, $8$ or $16$; in the
planar case a $4$-cycle already occurs. The engine is an extremal bound of
independent interest: a connected graph of minimum degree at least $2$ with no
$\Cfour$ and no $\Ceight$, $2$-cell embedded in a surface $S$, satisfies
$|E|\le 2|V|-2\chi(S)$. Specialised to the sphere this gives
$\mathrm{ex}_{\mathcal P}(n,\{\Cfour,\Ceight\})\le 2n-4$, which appears to be
the first planar Tur\'an bound for a \emph{set} of forbidden cycle lengths, and
which we sharpen to $20|E|\le 39|V|-78$ under minimum degree $3$. All bounds are
shown to be attained or nearly attained by explicit graphs.
\end{abstract}

\section{Introduction}

All graphs are finite and simple. We write $n=|V(G)|$, $m=|E(G)|$, and
$L(G)$ for the \emph{cycle spectrum} of $G$, the set of lengths of its cycles.
A \emph{power cycle} is a cycle of length $2^k$ with $k\ge2$.

\begin{conjecture}[Erd\H{o}s and Gy\'arf\'as, 1994]\label{conj:eg}
Every graph with $\delta(G)\ge3$ contains a power cycle.
\end{conjecture}

Conjecture~\ref{conj:eg} is open. Erd\H{o}s later wrote that he had become
convinced it is false, but no counterexample is known. The known proved cases
are: $K_{1,m}$-free graphs with $\delta\ge m+1$, and graphs with
$\Delta\ge2m-1$ (Shauger); planar claw-free graphs (Daniel and Shauger);
$3$-connected cubic planar graphs, which contain a $2^m$-cycle for some
$2\le m\le 7$ (Heckman and Krakovski); graphs with no long induced path (Gao and
Shan; Hu and Shen; Hegde, Sandeep and Shashank); and graphs of diameter $2$
(Carr). A theorem of Liu and Montgomery implies the conjecture for all graphs of
average degree exceeding an absolute constant, so the conjecture is entirely a
bounded-degree question. On the computational side, Markstr\"om verified it for
all cubic graphs on fewer than $29$ vertices; Exoo constructed a cubic graph on
$78$ vertices with no cycle of length $4$, $8$ or $16$, and one on $450$
vertices with no cycle of length $4$, $8$, $16$ or $32$.

Notably, \emph{no} case of Conjecture~\ref{conj:eg} with a minimum-degree
hypothesis above $3$ appears in the literature, and no case for planar graphs
beyond the claw-free and $3$-connected cubic ones. We supply both.

\begin{theorem}\label{thm:main-planar}
Every planar graph with $\delta(G)\ge4$ contains a cycle of length $4$ or $8$.
The same conclusion holds for every graph of Euler genus at most $1$ with
$\delta(G)\ge4$, and for every $2$-connected graph of Euler genus at most $2$
with $\delta(G)\ge4$.
\end{theorem}

\begin{theorem}\label{thm:main-bip}
Every bipartite graph with $\delta(G)\ge3$ that embeds in the torus contains a
cycle of length $4$, $8$ or $16$. In particular Conjecture~\ref{conj:eg} holds
for bipartite graphs of orientable genus at most $1$. Every \emph{planar}
bipartite graph with $\delta(G)\ge3$ already contains a $4$-cycle.
\end{theorem}

Both theorems rest on the following extremal bound, which we believe is of
independent interest.

\begin{theorem}\label{thm:density}
Let $G$ be a $2$-connected simple graph with $n\ge4$, containing no cycle of
length $4$ and no cycle of length $8$, and $2$-cell embedded in a surface $S$.
Then
\[
  m\;\le\;2n-2\chi(S).
\]
\end{theorem}

For the sphere this reads $m\le 2n-4$, i.e.
$\mathrm{ex}_{\mathcal P}(n,\{\Cfour,\Ceight\})\le 2n-4$. The known planar
Tur\'an bounds for single cycles are $\mathrm{ex}_{\mathcal P}(n,\Cfour)\le
\tfrac{15}{7}(n-2)$ (Dowden) and $\mathrm{ex}_{\mathcal P}(n,\Ceight)\le
\tfrac{69}{25}(n-2)$ (Bai, Liu, Nie and Zhang); our bound is smaller than
either, and coincides with the exact value $\mathrm{ex}_{\mathcal
P}(n,C_3)=2n-4$. We are not aware of any previous planar Tur\'an bound for a set
of two forbidden cycle lengths. Under the additional hypothesis
$\delta(G)\ge3$ we sharpen it to $20m\le 39n-78$ (Theorem~\ref{thm:sharp}).

The combinatorial heart of all of this is a short lemma with no topological
content whatsoever.

\begin{lemma}[Detour Lemma]\label{lem:detour}
Let $G$ contain no $4$-cycle. If some $5$-cycle of $G$ has three of its edges
lying in triangles, then $G$ contains an $8$-cycle.
\end{lemma}

\section{The Detour Lemma}

\begin{lemma}\label{lem:chordless}
In a $\Cfour$-free graph every $5$-cycle is chordless.
\end{lemma}

\begin{proof}
A chord of a $5$-cycle joins two vertices at distance $2$ along the cycle,
splitting it into arcs of lengths $2$ and $3$; the chord together with the arc of
length $3$ is a $4$-cycle.
\end{proof}

\begin{lemma}\label{lem:apex-off}
Let $G$ be $\Cfour$-free, let $C$ be a $5$-cycle of $G$, and let $uv\in E(C)$ lie
in a triangle $uvw$. Then $w\notin V(C)$, and $uv$ lies in exactly one triangle.
\end{lemma}

\begin{proof}
Write $C=uvx_1x_2x_3$. Suppose $w\in V(C)$, so $w\in\{x_1,x_2,x_3\}$. On a
$5$-cycle each of $x_1,x_2,x_3$ is adjacent along $C$ to at most one of $u,v$;
hence at least one of $wu$, $wv$ is a chord of $C$, contradicting
Lemma~\ref{lem:chordless}. If $uv$ lay in two triangles with distinct apexes
$w\ne w'$, then $w u w' v$ would be a $4$-cycle.
\end{proof}

\begin{lemma}\label{lem:apex-distinct}
Let $G$ be $\Cfour$-free and let $C$ be a $5$-cycle. If $e\ne e'$ are edges of
$C$ lying in triangles with apexes $w,w'\notin V(C)$, then $w\ne w'$.
\end{lemma}

\begin{proof}
Suppose $w=w'$ and write $C=v_1v_2v_3v_4v_5$. Up to the symmetry of $C$ we may
take $e=v_1v_2$ and $e'\in\{v_2v_3,\;v_3v_4,\;v_4v_5\}$.

If $e'=v_2v_3$ then $w$ is adjacent to $v_1,v_2,v_3$ and $w v_1 v_2 v_3$ is a
$4$-cycle. If $e'=v_3v_4$ then $w$ is adjacent to $v_1,v_2,v_3,v_4$ and
$w v_4 v_5 v_1$ is a $4$-cycle. If $e'=v_4v_5$ then $w$ is adjacent to
$v_1,v_2,v_4,v_5$ and $w v_2 v_3 v_4$ is a $4$-cycle. In each case the four
listed vertices are distinct, because $w\notin V(C)$.
\end{proof}

\begin{proof}[Proof of Lemma~\ref{lem:detour}]
Let $C$ be a $5$-cycle with edges $e_1,e_2,e_3$ lying in triangles with apexes
$w_1,w_2,w_3$. By Lemma~\ref{lem:apex-off} the $w_i$ lie off $C$, and by
Lemma~\ref{lem:apex-distinct} they are pairwise distinct. Replace each $e_i=uv$
by the path $u\,w_i\,v$. The resulting closed walk visits the five vertices of
$C$ in cyclic order with $w_1,w_2,w_3$ inserted, so it visits eight pairwise
distinct vertices, each once, and consecutive vertices are adjacent. Its length
is $5-3+3\cdot2=8$.
\end{proof}

\begin{remark}
Lemma~\ref{lem:detour} is equivalent to: in a $\{\Cfour,\Ceight\}$-free graph,
every $5$-cycle has at most two edges lying in triangles. We verified it by
exhaustive computation over all $143{,}038$ connected $\Cfour$-free graphs of
order at most $11$, of which $1{,}366$ contain a $5$-cycle with three edges in
triangles; every one of those contains an $8$-cycle.
\end{remark}

\section{The density theorem}

We use standard facts about $2$-cell embeddings. If $G$ is $2$-connected and
$2$-cell embedded in $S$, every face boundary is a cycle of $G$ and every edge
lies on exactly two distinct faces. Euler's formula reads $n-m+f=\chi(S)$, where
$f$ is the number of faces.

\begin{proof}[Proof of Theorem~\ref{thm:density}]
Fix a $2$-cell embedding. Since every face is bounded by a cycle and $G$ has no
$4$-cycle and no $8$-cycle, there are no faces of length $4$ and none of length
$8$. Assign initial charges
\[
  \mu(v)=d(v)-4\quad(v\in V),\qquad \mu(f)=\ell(f)-4\quad(f\text{ a face}).
\]
Since $\sum_v d(v)=2m$ and $\sum_f \ell(f)=2m$, and $n+f=m+\chi(S)$,
\[
  \sum_v\mu(v)+\sum_f\mu(f)=(2m-4n)+(2m-4f)=4m-4(n+f)=-4\chi(S).
\]

\emph{Discharging rule.} For every edge $e$ having a face of length $3$ on one
side and a face $g$ on the other, $g$ sends $\tfrac13$ to that $3$-face across
$e$.

The rule is well defined. Each edge lies on two distinct faces, and no two
$3$-faces are adjacent: if triangles $uvw$ and $uvw'$ bounded the two faces at
$uv$, then $w\ne w'$ (a face is determined by its boundary and its side) and
$w u w' v$ would be a $4$-cycle. Consequently a $3$-face never sends charge, and
each face $g$ sends at most $\tfrac13$ across each of its $\ell(g)$ edge-sides.

Let $\mu^{*}$ denote the final charges. Vertices neither send nor receive, so
$\sum_v\mu^{*}(v)=2m-4n$. For faces:

\begin{itemize}
\item $\ell(f)=3$: $f$ receives $\tfrac13$ across each of its three edges and
sends nothing, so $\mu^{*}(f)=-1+1=0$.
\item $\ell(f)=4$ or $\ell(f)=8$: no such face exists.
\item $\ell(f)=5$: the boundary is a $5$-cycle. If three of its edges had
$3$-faces on the other side, those three edges would lie in triangles and
Lemma~\ref{lem:detour} would produce an $8$-cycle. Hence $f$ sends at most
$2\cdot\tfrac13$, and $\mu^{*}(f)\ge 1-\tfrac23=\tfrac13>0$.
\item $\ell(f)=\ell\ge6$: $\mu^{*}(f)\ge \ell-4-\tfrac{\ell}{3}
=\tfrac{2\ell-12}{3}\ge0$.
\end{itemize}

Therefore $\sum_f\mu^{*}(f)\ge0$, and
\[
  2m-4n=\sum_v\mu^{*}(v)=-4\chi(S)-\sum_f\mu^{*}(f)\le-4\chi(S),
\]
which is the assertion.
\end{proof}

\begin{corollary}\label{cor:planarturan}
Every connected planar graph with $\delta\ge2$, $n\ge4$ and no $\Cfour$ and no
$\Ceight$ satisfies $m\le 2n-4$. In particular
$\mathrm{ex}_{\mathcal P}(n,\{\Cfour,\Ceight\})\le 2n-4$.
\end{corollary}

\begin{proof}
For $2$-connected $G$ this is Theorem~\ref{thm:density} with $\chi=2$, using that
every embedding of a connected graph in the sphere is $2$-cell. In general let
$B_1,\dots,B_b$ be the blocks, with $n_i,m_i$ vertices and edges. Then
$\sum_i n_i=n+b-1$ and $\sum_i m_i=m$. If $n_i\ge4$ then $B_i$ is $2$-connected
and $m_i\le 2n_i-4$; if $n_i=3$ then $B_i$ is a triangle and $m_i=2n_i-3$; if
$n_i=2$ then $m_i=1=2n_i-3$. In all cases $m_i\le 2n_i-3$, so
$m\le 2(n+b-1)-3b=2n-2-b$. For $b\ge2$ this gives $m\le 2n-4$, and for $b=1$ the
graph is $2$-connected.
\end{proof}

We verified Corollary~\ref{cor:planarturan} exhaustively over all $150{,}791$
connected $\Cfour$-free planar graphs with $\delta\ge2$ and $4\le n\le13$, of
which $8{,}242$ are also $\Ceight$-free.

\section{Minimum degree four}

\begin{proof}[Proof of Theorem~\ref{thm:main-planar}]
Suppose $G$ has $\delta(G)\ge4$, no $\Cfour$ and no $\Ceight$. Passing to a
component we may assume $G$ connected. Since $\delta\ge4$, no block of $G$ is a
single edge; let $B$ be an endblock, that is, a block containing at most one cut
vertex $c$ of $G$. Then $B$ is $2$-connected, every $v\in V(B)\setminus\{c\}$ has
$d_B(v)=d_G(v)\ge4$, and $d_B(c)\ge2$. Hence $|V(B)|\ge5$ and
\[
  2|E(B)|\;\ge\;4\bigl(|V(B)|-1\bigr)+2,\qquad\text{i.e.}\qquad
  |E(B)|\;\ge\;2|V(B)|-1 .
\]
$B$ inherits $\{\Cfour,\Ceight\}$-freeness and embeds in the same surfaces as
$G$.

\emph{Euler genus $0$ or $1$.} Embed $B$ in a surface $S$ of its own Euler
genus $\gamma\le1$. If $\gamma=0$ this is a sphere embedding of a connected
graph, hence $2$-cell. If $\gamma=1$ then $B$ is non-planar; a face of an
embedding in the projective plane that is not a disc is a M\"obius band, whose
complement is a disc, which would make $B$ planar; so the embedding is $2$-cell.
Now $\chi(S)=2-\gamma\ge1$ and Theorem~\ref{thm:density} gives
$|E(B)|\le 2|V(B)|-2$, contradicting the display above.

\emph{Euler genus $2$, $G$ $2$-connected.} Take $B=G$ and fix a minimum-genus
embedding, which is $2$-cell. We may assume the Euler genus is exactly $2$, so
$\chi=0$ and Theorem~\ref{thm:density} gives $m\le 2n$, while $\delta\ge4$ gives
$m\ge2n$. Hence $m=2n$ and every inequality in the proof of
Theorem~\ref{thm:density} is tight. Therefore $\sum_v(d(v)-4)=0$ with all
degrees at least $4$, so $G$ is $4$-regular; and $\mu^{*}(f)=0$ for every face,
so, since $\mu^{*}>0$ for faces of length $5$ and of length $\ge7$, and no
$4$-faces exist, every face has length $3$ or $6$, and every $6$-face sends the
full $6\cdot\tfrac13$, i.e.\ all six of its edge-sides border a $3$-face.

Every edge separates a $3$-face from a $6$-face: two $3$-faces on one edge is
excluded as before, and if both faces at an edge had length $6$ then some
$6$-face would have an edge-side not bordering a $3$-face.

Every $6$-face is bounded by a $6$-cycle. Its boundary walk does not backtrack
(that would force a vertex of degree $1$). If it is not a cycle it splits into
two closed walks of lengths at least $3$ summing to $6$, so it has the form
$w,p,q,w,r,s$ with $wpq$ a triangle. The edge $wp$ has the $6$-face on one side
and a $3$-face $T'$ on the other. If $T'$ is bounded by $wpq$ then all three
edges of $wpq$ separate these two faces, forcing $\deg(p)=2$, contradicting
$4$-regularity; otherwise $T'$ is bounded by $wpt$ with $t\ne q$, and the
triangles $wpq$, $wpt$ share $wp$, giving the $4$-cycle $q\,w\,t\,p$.

Since consecutive faces around a vertex share an edge and every edge is of type
$(3,6)$, the faces around each vertex alternate $3,6,3,6$. Fix $v$ with
$N(v)=\{a,b,c,d\}$ in rotation and faces $T_1$ (between $va,vb$), $H_1$ (between
$vb,vc$), $T_2$ (between $vc,vd$), $H_2$ (between $vd,va$), with
$\ell(T_i)=3$, $\ell(H_i)=6$. Then $T_1$ is the triangle $vab$, so $ab\in E$;
$T_2$ is the triangle $vcd$, so $cd\in E$; and $H_1$ is a $6$-cycle
$v\,b\,x\,y\,z\,c$. Consider the closed walk
\[
  W^{*}:\quad v,\;a,\;b,\;x,\;y,\;z,\;c,\;d,\;v,
\]
all of whose consecutive pairs are edges of $G$. Its eight vertices are
distinct: $a,b,c,d$ are distinct and differ from $v$; $b,x,y,z,c$ lie on the
$6$-cycle $H_1$ and differ from $v$; and
\begin{itemize}
\item $a=x$ would make $ab$ the edge $bx$ of $H_1$, so $T_1$ and $H_1$ would both
have $bv,ba$ as their edges at $b$, forcing $\deg(b)=2$;
\item $a=y$ would make $va$ a chord joining antipodal vertices of $H_1$,
splitting it into two $4$-cycles;
\item $a=z$ would give $a\sim c$, so $vac$ and $vab$ are triangles sharing $va$,
giving the $4$-cycle $b\,v\,c\,a$;
\item $d=z$, $d=y$, $d=x$ are excluded symmetrically.
\end{itemize}
Hence $W^{*}$ is an $8$-cycle, a contradiction.
\end{proof}

\begin{remark}[Sharpness]
The hypothesis $\delta\ge4$ cannot be weakened to $\delta\ge3$: the truncated
icosahedron (the buckyball) is a $3$-connected cubic planar graph on $60$
vertices with no $\Cfour$ and no $\Ceight$, its smallest power cycle being a
$\Csixteen$. The conclusion cannot be strengthened by dropping either forbidden
length: the icosidodecahedron is a $4$-regular planar $\Cfour$-free graph on
$30$ vertices, and it contains an $8$-cycle, as Theorem~\ref{thm:main-planar}
requires. Indeed the following holds.
\end{remark}

\begin{proposition}
Every planar graph with $\delta\ge4$ and no $\Cfour$ has at least $30$ vertices,
and $30$ is attained by the icosidodecahedron.
\end{proposition}

\begin{proof}
Take an endblock $B$ as above, so $|E(B)|\ge2|V(B)|-1$, and fix a plane
embedding. All faces are bounded by cycles, there is no $4$-face, and no two
$3$-faces share an edge, so each edge lies on at most one $3$-face and
$3f_3\le|E(B)|$. Hence
\[
  2|E(B)|=\sum_f\ell(f)\ \ge\ 3f_3+5(f-f_3)=5f-2f_3\ \ge\ 5f-\tfrac23|E(B)|,
\]
so $f\le\tfrac{8}{15}|E(B)|$, and Euler gives
$|V(B)|=2+|E(B)|-f\ge2+\tfrac{7}{15}|E(B)|$, i.e.\
$|E(B)|\le\tfrac{15}{7}(|V(B)|-2)$. With $|E(B)|\ge2|V(B)|-1$ this yields
$14|V(B)|-7\le15|V(B)|-30$, so $|V(B)|\ge23$. If $G$ is $2$-connected we may take
$B=G$ and use $m\ge2n$, giving $|V|\ge30$; otherwise $G$ has two endblocks
sharing at most one vertex, so $n\ge23+23-1\ge30$.
\end{proof}

\section{A sharper planar bound}

\begin{lemma}\label{lem:facelemma}
Let $G$ be a $2$-connected plane graph with $\delta(G)\ge3$, no $\Cfour$ and no
$\Ceight$. For a face $f$ let $t(f)$ be the number of edges of $\partial f$ whose
other face is a triangle. Then
\[
  \ell(f)=5\Rightarrow t(f)\le2,\qquad
  \ell(f)=6\Rightarrow t(f)\le1,\qquad
  \ell(f)=7\Rightarrow t(f)=0.
\]
\end{lemma}

\begin{proof}
Write $C=\partial f$, a cycle of length $\ell\le7$, and let $e_1,\dots,e_t$ be
the counted edges with triangle apexes $w_1,\dots,w_t$.

First, $w_i\notin V(C)$. For $\ell=5$ this is Lemma~\ref{lem:apex-off}. For
$\ell\in\{6,7\}$, write $C=u\,v\,x_1\cdots x_{\ell-2}$ with $e_i=uv$ and suppose
$w_i=x_j$. If $j=1$, then both $uv$ and $vx_1$ lie on $\partial f$ and on the
triangle $T=uvx_1$, so the two faces at each of these edges are $f$ and $T$; a
third edge at $v$ would have to be drawn inside $f$ or inside $T$, so
$\deg(v)=2$, contradicting $\delta\ge3$. The case $j=\ell-2$ is symmetric. For
$2\le j\le\ell-3$ put $a=j$ and $b=\ell-1-j$, the two lengths of the $C$-paths
from $v$ to $w_i$ and from $w_i$ to $u$; then $a,b\ge2$ and $a+b=\ell-1\le6$, so
$\min(a,b)\le3$. If $a=2$ then $u\,v\,x_1\,w_i$ is a $4$-cycle; if $a=3$ then
$v\,x_1\,x_2\,w_i$ is a $4$-cycle; and symmetrically for $b$.

Second, the $w_i$ are pairwise distinct. For $\ell=5$ this is
Lemma~\ref{lem:apex-distinct}. For $\ell\in\{6,7\}$ suppose $w_i=w_j=w$. If
$e_i,e_j$ share a vertex, say $e_i=tu$ and $e_j=uv$ with $t\ne v$, then
$t\,w\,v\,u$ is a $4$-cycle. Otherwise write $C=u\,v\,P\,s\,t\,Q\,u$ with
$e_i=uv$, $e_j=st$, $|P|=p\ge1$, $|Q|=q\ge1$ and $p+q+2=\ell\le7$, so
$\min(p,q)\le2$. If $p=1$ then $u\,w\,s\,v$ is a $4$-cycle; if $p=2$, say
$v\,x\,s$, then $v\,w\,s\,x$ is a $4$-cycle; symmetrically for $q$.

Now replace any subset of $\{e_1,\dots,e_t\}$ by the corresponding detours as in
Lemma~\ref{lem:detour}: this produces cycles of every length in
$[\ell,\ell+t]$. If $\ell=5$ and $t\ge3$, or $\ell=6$ and $t\ge2$, or $\ell=7$
and $t\ge1$, that interval contains $8$.
\end{proof}

\begin{theorem}\label{thm:sharp}
Let $G$ be a $2$-connected plane graph with $\delta(G)\ge3$, $n\ge4$, no
$\Cfour$ and no $\Ceight$, and let $f_3$ be its number of triangular faces. Then
\[
  7m\;\le\;13n+f_3-26,\qquad\text{and hence}\qquad 20m\;\le\;39n-78 .
\]
\end{theorem}

\begin{proof}
Run the discharging of Theorem~\ref{thm:density}. By Lemma~\ref{lem:facelemma}
the final charge of a face of length $5$ is at least $1-\tfrac23=\tfrac13$; of
length $6$ at least $2-\tfrac13=\tfrac53$; of length $7$ at least $3$; and of
length $\ell\ge9$ at least $(2\ell-12)/3\ge2$. There are no faces of length $4$
or $8$. Hence every non-triangular face retains at least $\tfrac13$, so
$\sum_f\mu^{*}(f)\ge (f-f_3)/3$. On the other hand
$\sum_f\mu^{*}(f)=\sum_f\mu(f)=-8-(2m-4n)=4n-2m-8$. With $f=m-n+2$ this gives
\[
  4n-2m-8\;\ge\;\tfrac13\bigl(m-n+2-f_3\bigr),
\]
i.e.\ $12n-6m-24\ge m-n+2-f_3$, i.e.\ $7m\le13n+f_3-26$. Finally,
$\Cfour$-freeness makes triangles pairwise edge-disjoint, so $3f_3\le m$, and
substituting gives $21m\le39n+m-78$.
\end{proof}

\begin{corollary}
Every $2$-connected planar graph with $\delta\ge3$ and $20m\ge39n-77$ contains a
cycle of length $4$ or $8$. Equivalently, writing $D=\sum_{d(v)\ge4}(d(v)-3)
=2m-3n$, this holds whenever $D\ge(9n-77)/10$. In particular, every
$2$-connected planar graph all of whose degrees are $3$ or $4$ and which has at
most $(n+77)/10$ vertices of degree $3$ contains a cycle of length $4$ or $8$.
\end{corollary}

Taking no vertices of degree $3$ recovers the planar case of
Theorem~\ref{thm:main-planar}, so this is a strict extension. The constant in
Theorem~\ref{thm:sharp} cannot be pushed below $90/59$: chaining $k$ copies of
the buckyball at cut vertices produces planar $\{\Cfour,\Ceight\}$-free graphs
with $\delta\ge3$, $n=59k+1$ and $m=90k$.

\begin{remark}
The same discharging yields a further structural restriction which may be useful
elsewhere. If $F$ is a set of faces of a $2$-connected plane graph whose closed
union is a disc, then $\partial F$ is a cycle of length
$\sum_{f\in F}\ell(f)-2e_F$, where $e_F$ counts edges incident with two faces of
$F$. Applying this to two faces sharing exactly one edge shows that in a
$\{\Cfour,\Ceight\}$-free plane graph no two $3$-faces are adjacent, \emph{no two
$5$-faces are adjacent}, and no $3$-face is adjacent to a $7$-face.
\end{remark}

\section{Bipartite graphs of genus at most one}

\begin{proof}[Proof of Theorem~\ref{thm:main-bip}]
Let $G$ be bipartite with $\delta(G)\ge3$, embeddable in the torus, and suppose
$G$ has no $\Cfour$; being bipartite, $G$ then has girth at least $6$. Passing
to a component and then to an endblock $B$ with cut vertex $c$ (or $B=G$), we
have $B$ $2$-connected with $d_B(v)=d_G(v)\ge3$ for $v\ne c$ and $d_B(c)\ge2$,
so
\[
  2m\;\ge\;3n-1,
\]
where from now on $n=|V(B)|$ and $m=|E(B)|$. Fix a minimum-genus embedding of
$B$, which is $2$-cell. Since $B$ is $2$-connected every face is bounded by a
cycle, necessarily of even length at least $6$; hence $2m=\sum_f\ell(f)\ge6f$,
i.e.\ $f\le m/3$.

\emph{Genus $0$.} Euler gives $f=m-n+2\le m/3$, so $m\le\tfrac32(n-2)$ and
$2m\le3n-6<3n-1$, a contradiction. Hence every planar bipartite graph with
$\delta\ge3$ contains a $4$-cycle.

\emph{Genus $1$.} Euler gives $f=m-n\le m/3$, so $2m\le3n$, and therefore
$3n-1\le2m\le3n$. If $n$ is odd then $3n$ is odd, so $2m=3n-1$ and
$f=(n-1)/2$; then $\sum_f\ell(f)=2m=3n-1$ while $6f=3n-3$, and since all face
lengths are even and at least $6$, exactly one face has length $8$, giving an
$8$-cycle. If $n$ is even then $2m=3n$, so $B$ is cubic, $f=n/2$ and
$\sum_f\ell(f)=6f$: every face is a hexagon.

So $B$ is a hexagonal tessellation of the torus. Its universal cover is the
hexagonal tessellation $\mathcal H$ of the plane, and the deck group is a rank-2
lattice $\Lambda'$ of translations, so $B=\mathcal H/\Lambda'$. Identify the
translation lattice of $\mathcal H$ with $\mathbb Z^2$ and let $|\cdot|$ be the
triangular norm, $|(a,b)|=|a|+|b|$ if $ab\ge0$ and $\max(|a|,|b|)$ otherwise;
two vertices of $\mathcal H$ have the same image in $B$ exactly when they have
the same type and their difference lies in $\Lambda'$. Put
\[
  t=\min\{\,|\lambda| : \lambda\in\Lambda'\setminus\{0\}\,\}.
\]
Then $t\le1$ makes the quotient non-simple; $t=2$ produces a $4$-cycle; $t=3$
produces an $8$-cycle; and $t\ge4$ produces a $16$-cycle. For the last, take the
polyhex whose four hexagon centres are $0$, $e_1$, $2e_1$, $e_2$: it has four
hexagons and four adjacent pairs, so its boundary has length
$6\cdot4-2\cdot4=16$, and every difference of two same-type vertices along that
boundary has norm at most $3$, so the projection to $B$ is injective when
$t\ge4$. For $t\in\{2,3,4\}$ and each of the $12$, $18$, respectively $24$
lattice vectors $\lambda$ of that norm we exhibit a path of length $2t$ in
$\mathcal H$ from a vertex to its $\lambda$-translate all of whose same-type
differences have norm smaller than $t$; its projection is then a $2t$-cycle. This
is a finite verification which we carried out exhaustively, with no exceptions.
\end{proof}

\begin{remark}[Sharpness]
All three lengths in Theorem~\ref{thm:main-bip} are needed. The hexagonal torus
$\mathcal H/\Lambda'$ with $t=5$ of smallest index, on $38$ vertices, is a cubic
bipartite graph of girth exactly $6$ with cycle counts $19,0,171,323,969$ at
lengths $6,8,10,12,14$: it has neither a $4$-cycle nor an $8$-cycle, and its
smallest power cycle is a $16$-cycle. Since $t\ge5$ forces the triangular balls
of radius $2$ about lattice points to be disjoint, and each contains $19$ points,
the index of $\Lambda'$ is at least $19$; hence $38$ is exactly the minimum order
of a toroidal bipartite graph with $\delta\ge3$, girth $6$ and no $8$-cycle. In
particular the implication ``bipartite, $\delta\ge3$, girth $6$ $\Rightarrow$
$\Ceight$'' is false.
\end{remark}

\begin{remark}
The Klein bottle and the projective plane are not covered by
Theorem~\ref{thm:main-bip}: the counting above already forces Euler genus at
least $2$ in the relevant range, and the classification of hexagonal
tessellations we use is specific to the torus.
\end{remark}

\section{Concluding remarks}

Theorem~\ref{thm:density} degenerates exactly at $\chi=0$, where the bound
$m\le 2n$ meets the hypothesis $m\ge2n$ supplied by $\delta\ge4$; this is why
the torus and Klein bottle require the separate argument in the proof of
Theorem~\ref{thm:main-planar}, and why surfaces of negative Euler
characteristic are out of reach by this method. For $\delta\ge5$ the same
counting gives $n\le-4\chi(S)$, so for each fixed surface only finitely many
orders survive.

The planar case of Conjecture~\ref{conj:eg} at $\delta=3$ remains open, and two
obstructions delimit it. First, one cannot hope to prove that a planar
counterexample is cubic: identifying a vertex of one buckyball with a vertex of
another yields a planar graph with $\delta\ge3$, a vertex of degree $6$, and no
$\Cfour$ and no $\Ceight$. Second, one cannot hope to conclude with the lengths
$4$, $8$ and $16$ alone: the buckyball is triangle-free and the truncated
dodecahedron is pentagon-free, both are $3$-connected cubic planar with no
$\Cfour$ and no $\Ceight$, and every $2$-connected plane graph with $\delta\ge3$
and no $4$-face satisfies $\sum_f(6-\ell(f))\ge12$ and therefore has a $3$-face
or a $5$-face. This is consistent with the range $2\le m\le7$ in the theorem of
Heckman and Krakovski, and with Exoo's $3$-connected cubic planar graph on $420$
vertices having no cycle of length $2^m$ for $m\le4$.

\begin{question}
Is $\mathrm{ex}_{\mathcal P}(n,\{\Cfour,\Ceight\})=\bigl(\tfrac32+o(1)\bigr)n$?
The constructions above give $\tfrac{90}{59}n+O(1)$ from below, and
Theorem~\ref{thm:sharp} gives $\tfrac{39}{20}n$ from above under $\delta\ge3$.
\end{question}

\end{document}
